% SPDX-FileCopyrightText: 2026 Hasan Melih Akbulut
% SPDX-License-Identifier: CC-BY-4.0

\chapter{What exists, what does not, and what comes next}
\label{ch:conclusion}

\section{What exists}

A radiation-hardened neuromorphic system-on-chip for the IHP SG13G2
130\,nm open process design kit, in register-transfer form and in
layout, built end to end with open tools.

In register-transfer form: a 32-bit RISC-V application core with a
single-error-correcting register file and a scrub engine that walks it;
a spiking accelerator with its event engine, its address-event queues
and its serial attachment to the core; a bus fabric with a real
peripheral bus behind it; a time base, a general-purpose timer and a
three-stage watchdog; memories behind error-correcting codes; a boot
path that is hardened rather than assumed; and triple-modular
redundancy where the design decided it was worth its area.

In layout: a placed and routed die of $2306.40 \times
2075.22$\,\si{\micro\metre} carrying eight vendor memory macros and
168{,}592 instances, with zero routing design-rule errors, a layout
versus schematic comparison that matches uniquely, a zero-difference
geometric comparison against the extracted view, and positive hold
slack at all three corners. Every geometric error the sign-off decks do
report lies inside a vendor macro or within half a micrometre of one:
the design is clean where the design drew.

And, throughout, an instrumentation layer that is part of the result
rather than scaffolding around it: guards that count redundancy in the
netlist because the source cannot be trusted to survive synthesis,
formal property sets that state one-line invariants over all reachable
states, fault-injection campaigns at both register-transfer and gate
level, and a cross-check that asks whether an any-state theorem and a
campaign agree about the same block.

\section{What does not exist, stated plainly}

No silicon of this system-on-chip. A frozen slice of the accelerator
was prepared for a multi-project shuttle and is pinned by hash; the
submission itself is an owner decision with a deadline and a cost, and
nothing in this document should be read as saying the part has been
fabricated.

No frequency. The project adopted a rule early --- report periods,
slacks and corner names, never a megahertz figure --- and has kept it
through every layout. A reader who wants a speed grade will not find
one here, and that is deliberate.

No mesh. The multi-node link that would make this a neuromorphic
\emph{fabric} rather than a single accelerator is specified down to its
24-bit link word and its per-node counters, and no register-transfer
code implements any of it. By the project's own ranking it is the
highest-leverage unbuilt item in the design.

No closed multiplier proof. Six instruction checks in the formal
programme have no verdict on either core, under two solvers, at a
stated bound of two hours each. The absence is bounded and named; it is
not a pass.

No verdict for the register-file consistency check on the substituted
core at its required depth. What exists instead is a stronger statement
about a smaller thing --- the register file and its scrub proved
unbounded with the codec abstracted --- and a located reason for the
original difficulty.

No flight heritage, no qualification campaign, and no claim that any of
the hardening in this design has been validated by irradiation. Every
hardening claim in this document is a claim about a mechanism, its
cost, and what an injected-fault campaign observed in simulation.

\section{The open items, ranked by what they would settle}

\begin{enumerate}
\item \textbf{The mesh link and a second node.} The largest unbuilt
      thing, and the one that changes what the part is rather than how
      well it does what it already does.
\item \textbf{A registered request phase in the fabric.} The
      register-file read into the peripheral space is the dominant
      violating path in every layout this project has produced, and no
      placement experiment has moved it, because it is a structure and
      not a placement. It is the one change that would change the worst
      number in \cref{tab:headline}, and it is unpriced.
\item \textbf{Hold margin as a first-class budget.} Three separate
      improvements have now been refused because they cost more hold
      than the design has. The next layout study worth running is not
      another improvement but a hold-repair pass aimed at the fast
      corner, to find out whether the 29.6\,ps is structural or
      tunable.
\item \textbf{Irradiation.} Everything in \cref{ch:harden} is a
      mechanism with a simulated campaign behind it. The step that
      would turn it into evidence about silicon is a beam, and it is
      not a step this project can take alone.
\item \textbf{The remaining peripherals.} A second serial port, the
      synchronous and two-wire interfaces, and the space-oriented link
      and bus codecs are specified and unbuilt; one of them is now
      unblocked by a bridge that exists and is proved, and none of them
      is on the critical path of anything above.
\end{enumerate}

\section{What this document is}

A record, not a report. Its distinguishing habit is the one stated in
the introduction and applied in every chapter since: a measurement that
is superseded is marked with its date and left standing, because the
history of a number is evidence about the number. A check that did not
finish is reported in the word its tool used. A result whose scope is
narrower than its headline says so in the same paragraph as the
headline. A run that measured the wrong thing is kept and labelled
rather than deleted, because the failure shape is more instructive than
the result would have been.

That habit has a cost --- this document is longer and less quotable
than a datasheet --- and one benefit that seems worth it: everything in
it can be checked, and the places where it is weakest are the places it
says so itself.
